\chapter[1996 --- Doherty et al.]{1996: Peter C. Doherty, Rolf M. Zinkernagel}
\label{ch:1996_Doherty}

\section*{Main Contribution}

\textit{Discovered how T cells recognize virus-infected cells, explaining the specificity of the cell-mediated immune response and how the immune system targets infected cells.}\cite{nobel-1996,lecture-1996}

\section{Official Citation}

for their discoveries concerning the specificity of the cell mediated immune defence

\section{Background and Historical Context}

The immune system's ability to distinguish self from non-self is one of the most notable feats of biological recognition. While the humoral arm of the immune system---mediated by antibodies---had been extensively studied since the late nineteenth century, the cellular arm---mediated by T lymphocytes---remained poorly understood until the 1970s. The discovery by Peter C. Doherty and Rolf M. Zinkernagel that cytotoxic T cells recognize foreign antigens only when presented on the surface of infected cells in association with major histocompatibility complex (MHC) molecules was a paradigm-shifting insight that earned them the 1996 Nobel Prize in Physiology or Medicine and fundamentally transformed our understanding of cellular immunity.

Peter Colin Doherty was born on October 15, 1940, in Brisbane, Australia. He studied veterinary science at the University of Queensland, earning his B.V.Sc.\ in 1962, and then completed his Ph.D.\ in immunology at the Australian National University in Canberra in 1970 under the supervision of Gordon Ada. Doherty's doctoral work concerned the immune response to viral infections in mice, and it was during this period that he began the collaboration with Zinkernagel that would lead to their Nobel Prize--winning discovery.

Rolf Martin Zinkernagel was born on January 4, 1944, in Riehen, Switzerland. He studied medicine at the University of Basel, earning his M.D.\ in 1968, and then completed his doctoral work at the Australian National University in Canberra under the supervision of George Mackaness. Zinkernagel's training in cellular immunology and viral pathogenesis provided the foundation for his collaborative work with Doherty.

The intellectual context of their work was shaped by the growing recognition that T lymphocytes---rather than antibodies---were the primary mediators of antiviral immunity and transplant rejection. In the 1960s and early 1970s, researchers had established that T cells could kill virus-infected cells directly, and that this killing was highly specific: T cells from a mouse infected with one virus would kill only cells infected with that virus, not cells infected with a different virus. The molecular basis of this specificity was unknown.

A parallel line of research concerned the major histocompatibility complex (MHC), a set of cell surface glycoproteins encoded by a highly polymorphic gene cluster. MHC molecules had been identified as the primary targets of the immune response in transplant rejection---hence their name---but their normal physiological function was unclear. The MHC in humans is called the human leukocyte antigen (HLA) system, and in mice it is called the H-2 complex.

\section{The Discovery/Contribution}

\subsection{The Discovery of MHC Restriction}

In 1974, Doherty and Zinkernagel published a landmark paper in \textit{Nature} that reported a startling observation. They were studying the immune response of mice to lymphocytic choriomeningitis virus (LCMV), a virus that infects the brain and causes inflammation. They measured the ability of cytotoxic T cells from infected mice to kill virus-infected target cells in vitro.

Their key finding was that cytotoxic T cells could only kill virus-infected target cells if the target cells shared the same MHC genotype as the mouse in which the T cells had been generated. In other words, T cells from an H-2$^k$ mouse would kill virus-infected H-2$^k$ cells but not virus-infected H-2$^b$ cells, even though both target cells were infected with the same virus. Conversely, T cells from an H-2$^b$ mouse would kill virus-infected H-2$^b$ cells but not H-2$^k$ cells.

This phenomenon, which they termed ``MHC restriction,'' was a significant discovery. It meant that the T cell receptor did not recognize the viral antigen alone---it recognized a composite structure consisting of the viral antigen bound to an MHC molecule on the surface of the infected cell. The T cell was thus ``restricted'' by the MHC: it could only ``see'' the antigen when it was presented in the context of a compatible MHC molecule.

\subsection{The Dual Recognition Model}

Doherty and Zinkernagel proposed that the T cell receptor recognized two molecular features simultaneously: the foreign antigen (a viral peptide fragment) and the self-MHC molecule. This ``dual recognition'' model explained both the specificity of the immune response (why T cells killed only virus-infected cells) and its MHC restriction (why T cells killed only cells with the same MHC type).

The dual recognition model was subsequently validated and refined by the identification of the T cell receptor as a heterodimeric protein composed of $\alpha$ and $\beta$ chains, each containing a variable region that recognizes the antigen--MHC complex. The crystal structure of the T cell receptor bound to an antigen--MHC complex, solved in the 1990s, confirmed the dual recognition model at atomic resolution.

\subsection{Antigen Processing and Presentation}

The discovery of MHC restriction also raised the question of how viral antigens were processed and presented on the cell surface. Doherty, Zinkernagel, and their colleagues showed that viral proteins synthesized inside infected cells were degraded into peptide fragments by a specialized proteolytic system called the proteasome. These peptide fragments were then transported into the endoplasmic reticulum by a transporter associated with antigen processing (TAP) and loaded onto MHC class I molecules for presentation on the cell surface.

This antigen processing and presentation pathway---now known as the MHC class I pathway---is essential for the detection of intracellular pathogens by cytotoxic T cells. A parallel pathway, the MHC class II pathway, presents antigens from extracellular pathogens to helper T cells, orchestrating the broader immune response.

\section{Impact on Modern Medicine}

The discovery of MHC restriction by Doherty and Zinkernagel has had a significant impact on immunology and medicine.

\subsection{Vaccine Development}

Understanding how T cells recognize viral antigens in the context of MHC molecules has been essential for the development of vaccines, particularly those designed to elicit strong cellular immune responses. Modern vaccine strategies, including DNA vaccines, viral vector vaccines, and mRNA vaccines (such as those developed for COVID-19), are designed to deliver antigens into cells where they can be processed and presented on MHC class I molecules to stimulate cytotoxic T cell responses. The understanding that MHC restriction determines which peptides are presented has guided the selection of immunodominant epitopes for vaccine design.

\subsection{Autoimmune Disease}

The principle of MHC restriction has also illuminated the mechanisms of autoimmune disease. Many autoimmune diseases are strongly associated with specific MHC alleles: type 1 diabetes with HLA-DR3 and HLA-DR4, rheumatoid arthritis with HLA-DR4, and ankylosing spondylitis with HLA-B27. Understanding how specific MHC molecules present self-antigens to T cells has led to insights into the loss of self-tolerance and the development of autoimmune tissue damage.

\subsection{Transplant Medicine}

MHC restriction is also directly relevant to organ transplantation. The matching of donor and recipient MHC alleles (HLA typing) is a critical component of transplant compatibility, and the mechanisms of graft rejection involve the recognition of donor MHC molecules by recipient T cells. Advances in immunosuppressive therapy and in the understanding of transplant tolerance are informed by the principles of MHC restriction.

\section{Legacy}

Peter Doherty continued his research at the University of Melbourne and at St. Jude Children's Research Hospital in Memphis, Tennessee, where he studied the mechanisms of immunological memory and the immune response to influenza virus. He has been a prominent public communicator of science, authoring several popular science books and receiving the Australian of the Year award in 1997.

Rolf Zinkernagel became director of the Institute of Experimental Immunology at the University of Z{\"u}rich and continued to study the mechanisms of T cell--mediated immunity, with a particular focus on immunological memory and the requirements for long-lived protective immunity. He received numerous honors, including the Cloëtta Prize and the Rolf Schock Prize.

Together, Doherty and Zinkernagel revealed the molecular logic of cellular immunity---how T cells recognize and respond to intracellular pathogens. Their discovery of MHC restriction is one of the foundational principles of modern immunology, and its implications extend to vaccine development, autoimmune disease, transplant medicine, and cancer immunotherapy.


\section{Modern Developments}

Doherty and Zinkernagel's T cell work has led to modern immunotherapy. Understanding of T cell recognition has enabled development of checkpoint inhibitors (Nobel 2018). Understanding of viral epitopes has informed vaccine design, including universal influenza vaccines. Understanding of T cell exhaustion has led to combination immunotherapies for cancer. Understanding of cross-presentation has informed development of cancer vaccines and dendritic cell therapies.


\section{Bibliography}

\begin{thebibliography}{9}
\bibitem{nobel-1996}
Nobel Prize Outreach, ``The Nobel Prize in Physiology or Medicine 1996,''\emph{NobelPrize.org}.\\
\url{https://www.nobelprize.org/prizes/medicine/1996/summary/}

\bibitem{lecture-1996}
Peter C. Doherty, ``Nobel Lecture,''\emph{NobelPrize.org}. Winner-authored primary source; downloadable from the official Nobel Prize archive.\\
\url{https://www.nobelprize.org/prizes/medicine/1996/doherty/lecture/}
\end{thebibliography}
